\documentclass[11pt]{article}

\usepackage{amsmath,amssymb,amsfonts,bm}
\usepackage[margin=1in]{geometry}
\usepackage{mathtools}
\usepackage{amsthm}
\usepackage{enumitem}

\title{Dynamic 3D Gaussian Splatting with Explicit Real-Valued ODE-Based Mean and Covariance Evolution}
\date{}

\begin{document}

\maketitle

\section{Overview}

We propose a dynamic 3D Gaussian Splatting framework that preserves the original 3DGS rendering pipeline while replacing neural deformation fields with explicit ordinary differential equations. Each primitive remains a standard 3D anisotropic Gaussian in world space, and rendering is still performed exactly as in the original 3DGS method: each 3D Gaussian is projected onto the image plane, converted into a 2D Gaussian footprint, depth-sorted, and rasterized by alpha compositing.

The rendering backbone is unchanged. The proposed method does not introduce volumetric ray tracing, radiance-field integration, implicit neural warp fields, or MLP-based deformation modules. The only modification is that the Gaussian mean and covariance are modeled as explicit continuous-time functions generated by real-valued ordinary differential equations.

Unlike planar complex-valued formulations, the proposed framework models Gaussian motion directly in three-dimensional space. Each Gaussian center follows a real-valued ODE trajectory in $\mathbb{R}^3$, making the motion representation naturally compatible with the 3D Gaussian primitive itself. This creates a unified formulation in which both the mean and covariance of each Gaussian evolve directly in 3D world space.

The purpose of this framework is to study whether explicit ODE-based temporal evolution of Gaussian mean and covariance can represent deformable motion more effectively than common alternatives such as linear approximation or deformation-network-based methods. Since the renderer remains the original 3DGS renderer, the comparison isolates the effect of temporal modeling alone.

\section{Dynamic 3D Gaussian Representation}

Let the dynamic scene at normalized time $\tau$ be represented by a set of $N$ anisotropic Gaussian primitives
\begin{equation}
\mathcal{G}_{\tau}=\{G_i(\tau)\}_{i=1}^N,
\end{equation}
where primitive $i$ is defined as
\begin{equation}
G_i(\tau)=\bigl(\mu_i(\tau),\Sigma_i(\tau),\alpha_i,c_i\bigr).
\end{equation}
Here,
\begin{itemize}[leftmargin=1.5em]
    \item $\mu_i(\tau)\in\mathbb{R}^3$ is the time-dependent Gaussian mean,
    \item $\Sigma_i(\tau)\in\mathbb{R}^{3\times 3}$ is the time-dependent covariance matrix,
    \item $\alpha_i\in[0,1]$ is the Gaussian opacity,
    \item $c_i$ denotes the appearance parameter, such as RGB color or spherical harmonics coefficients.
\end{itemize}

To isolate the effect of temporal geometry modeling, we focus on the time evolution of the Gaussian mean and covariance. In the simplest setting, $\alpha_i$ and $c_i$ are kept fixed over time and handled exactly as in standard 3DGS.

We normalize the original timestamp $t$ into a bounded temporal coordinate $\tau\in[-1,1]$:
\begin{equation}
\tau = 2\frac{t-t_{\min}}{t_{\max}-t_{\min}}-1.
\end{equation}

\section{Canonical Gaussian State}

For each Gaussian primitive $i$, we define a canonical state at the reference time $\tau=0$:
\begin{equation}
G_i^0=\bigl(\mu_i^0,\Sigma_i^0,\alpha_i,c_i\bigr),
\end{equation}
where
\begin{itemize}[leftmargin=1.5em]
    \item $\mu_i^0\in\mathbb{R}^3$ is the canonical mean,
    \item $\Sigma_i^0\in\mathbb{R}^{3\times 3}$ is the canonical covariance,
    \item $\alpha_i$ is the opacity,
    \item $c_i$ is the appearance parameter.
\end{itemize}

The dynamic Gaussian at time $\tau$ is obtained by evolving the canonical mean and covariance directly through explicit ordinary differential equations starting from $(\mu_i^0,\Sigma_i^0)$. Thus, instead of assigning an unrelated Gaussian to every frame, each primitive follows a continuous-time trajectory in both position and shape. The temporal evolution is governed directly by ODE parameters, not by a warp field, implicit function, or MLP deformation module.

\section{Real-Valued 3D ODE for Mean Evolution}

The Gaussian mean is modeled directly as a three-dimensional time-dependent quantity. For each Gaussian primitive $i$, we define a displacement state
\begin{equation}
x_i(\tau)\in\mathbb{R}^3,
\end{equation}
where $x_i(\tau)$ represents the 3D displacement of the Gaussian center from its canonical position.

The time-dependent Gaussian mean is then given by
\begin{equation}
\mu_i(\tau)=\mu_i^0+x_i(\tau).
\end{equation}

The displacement state evolves according to a real-valued linear ordinary differential equation in 3D:
\begin{equation}
\frac{d x_i(\tau)}{d\tau}
=
A_i x_i(\tau)+b_i,
\end{equation}
where
\begin{equation}
A_i\in\mathbb{R}^{3\times 3},
\qquad
b_i\in\mathbb{R}^3
\end{equation}
are trainable ODE parameters for Gaussian $i$.

The initial condition is
\begin{equation}
x_i(0)=x_i^0,
\end{equation}
where $x_i^0\in\mathbb{R}^3$ is the learnable initial displacement at the canonical time. In the simplest implementation, one may initialize $x_i^0=0$, so that
\begin{equation}
\mu_i(0)=\mu_i^0.
\end{equation}

Therefore, the complete mean trajectory is
\begin{equation}
\mu_i(\tau)=\mu_i^0+x_i(\tau),
\qquad
\frac{d x_i(\tau)}{d\tau}=A_i x_i(\tau)+b_i.
\end{equation}

This formulation is fully three-dimensional. Unlike a planar representation, no local 2D basis or complex-valued state is required. The ODE state, Gaussian mean, and Gaussian primitive all live naturally in $\mathbb{R}^3$.

\subsection{Closed-Form Solution}

For the linear ODE
\begin{equation}
\frac{d x_i(\tau)}{d\tau}
=
A_i x_i(\tau)+b_i,
\end{equation}
the solution can be written as
\begin{equation}
x_i(\tau)
=
e^{A_i\tau}x_i^0
+
\int_0^\tau e^{A_i(\tau-s)}b_i\,ds.
\end{equation}

If $A_i$ is invertible, this becomes
\begin{equation}
x_i(\tau)
=
e^{A_i\tau}x_i^0
+
A_i^{-1}\left(e^{A_i\tau}-I\right)b_i.
\end{equation}

Thus, the Gaussian mean is
\begin{equation}
\mu_i(\tau)
=
\mu_i^0
+
e^{A_i\tau}x_i^0
+
A_i^{-1}\left(e^{A_i\tau}-I\right)b_i.
\end{equation}

When $A_i=0$, the model reduces to simple linear drift:
\begin{equation}
x_i(\tau)=x_i^0+b_i\tau,
\end{equation}
and therefore
\begin{equation}
\mu_i(\tau)=\mu_i^0+x_i^0+b_i\tau.
\end{equation}

This shows that linear motion is a special case of the proposed ODE formulation, while nonzero $A_i$ allows curved, coupled, accelerating, rotating, expanding, or contracting 3D trajectories.

\section{Explicit ODE-Based Covariance Evolution}

The covariance must remain symmetric positive definite for all $\tau$. To guarantee this, we use a rotation-scale decomposition
\begin{equation}
\Sigma_i(\tau)
=
R_i(\tau)
\Lambda_i(\tau)
R_i(\tau)^\top,
\end{equation}
where
\begin{equation}
R_i(\tau)\in SO(3)
\end{equation}
is a valid 3D rotation matrix and $\Lambda_i(\tau)$ is a positive diagonal scale matrix.

We define
\begin{equation}
\Lambda_i(\tau)
=
\operatorname{diag}
\left(
e^{s_{i1}(\tau)},
e^{s_{i2}(\tau)},
e^{s_{i3}(\tau)}
\right),
\end{equation}
which guarantees positive principal scales by construction.

\subsection{Covariance Scale Evolution}

The logarithmic scale variables evolve according to
\begin{equation}
\frac{d s_{ij}(\tau)}{d\tau}
=
\kappa_{ij},
\qquad
j\in\{1,2,3\},
\end{equation}
with initial conditions
\begin{equation}
s_{ij}(0)=s_{ij}^0.
\end{equation}

The closed-form solution is
\begin{equation}
s_{ij}(\tau)=s_{ij}^0+\kappa_{ij}\tau.
\end{equation}

Thus,
\begin{equation}
\Lambda_i(\tau)
=
\operatorname{diag}
\left(
e^{s_{i1}^0+\kappa_{i1}\tau},
e^{s_{i2}^0+\kappa_{i2}\tau},
e^{s_{i3}^0+\kappa_{i3}\tau}
\right).
\end{equation}

\subsection{Covariance Orientation Evolution}

To model 3D covariance orientation naturally, we define an angular velocity vector
\begin{equation}
\omega_i^\Sigma\in\mathbb{R}^3.
\end{equation}

Let $\widehat{\omega}_i^\Sigma\in\mathfrak{so}(3)$ be the skew-symmetric matrix associated with $\omega_i^\Sigma$:
\begin{equation}
\widehat{\omega}_i^\Sigma
=
\begin{bmatrix}
0 & -\omega_{iz}^\Sigma & \omega_{iy}^\Sigma\\
\omega_{iz}^\Sigma & 0 & -\omega_{ix}^\Sigma\\
-\omega_{iy}^\Sigma & \omega_{ix}^\Sigma & 0
\end{bmatrix}.
\end{equation}

The covariance rotation evolves according to the real-valued matrix ODE
\begin{equation}
\frac{dR_i(\tau)}{d\tau}
=
\widehat{\omega}_i^\Sigma R_i(\tau),
\end{equation}
with initial condition
\begin{equation}
R_i(0)=R_i^0.
\end{equation}

The solution is
\begin{equation}
R_i(\tau)
=
\exp\left(\widehat{\omega}_i^\Sigma \tau\right)R_i^0.
\end{equation}

This guarantees that $R_i(\tau)$ remains a valid 3D rotation matrix for all $\tau$.

Therefore, the covariance at time $\tau$ is
\begin{equation}
\Sigma_i(\tau)
=
R_i(\tau)
\Lambda_i(\tau)
R_i(\tau)^\top.
\end{equation}

This gives a unified 3D ODE-based description of Gaussian shape evolution: the mean follows a real-valued ODE in $\mathbb{R}^3$, the scale follows real-valued scalar ODEs, and the orientation follows a real-valued rotation ODE in $SO(3)$.

\section{Dynamic Gaussian at Time \texorpdfstring{$\tau$}{tau}}

At any time $\tau$, the ODE system produces the Gaussian mean and covariance
\begin{equation}
\mu_i(\tau),
\qquad
\Sigma_i(\tau),
\end{equation}
for each primitive $i$. The resulting dynamic Gaussian set is
\begin{equation}
\mathcal{G}_{\tau}
=
\left\{
\bigl(\mu_i(\tau),\Sigma_i(\tau),\alpha_i,c_i\bigr)
\right\}_{i=1}^N.
\end{equation}

Therefore, the entire temporal variation of the scene comes only from the ODE-controlled evolution of Gaussian mean and covariance.

\section{Standard 3DGS Rendering}

The dynamic Gaussian set at time $\tau$ is rendered using the original 3D Gaussian Splatting pipeline. Each 3D Gaussian is projected into the camera image plane, producing a 2D Gaussian footprint with projected mean and covariance
\begin{equation}
\mu_i'(\tau)\in\mathbb{R}^2,
\qquad
\Sigma_i'(\tau)\in\mathbb{R}^{2\times 2}.
\end{equation}

For a pixel coordinate $x\in\mathbb{R}^2$, the screen-space contribution of Gaussian $i$ is
\begin{equation}
g_i(x,\tau)
=
\exp
\left(
-\frac{1}{2}
(x-\mu_i'(\tau))^\top
\Sigma_i'(\tau)^{-1}
(x-\mu_i'(\tau))
\right).
\end{equation}

As in standard 3DGS, projected Gaussians are depth-sorted and alpha-composited. If $\pi(x)$ denotes the ordered set of relevant Gaussians for pixel $x$, the rendered color is
\begin{equation}
C(x,\tau)
=
\sum_{i\in\pi(x)}
T_i(x,\tau)\,
\alpha_i(x,\tau)\,
c_i,
\end{equation}
where
\begin{equation}
\alpha_i(x,\tau)
=
\alpha_i g_i(x,\tau),
\end{equation}
and
\begin{equation}
T_i(x,\tau)
=
\prod_{j<i}
\left(
1-\alpha_j(x,\tau)
\right).
\end{equation}

The rendering backbone is identical to the original 3DGS formulation. The only change is the way $\mu_i(\tau)$ and $\Sigma_i(\tau)$ are generated over time.

\section{Training Objective}

The model is trained end-to-end from dynamic image sequences. For each observed image $I_t$ at time $t$, we convert the timestamp to $\tau$, solve the explicit ODEs to obtain Gaussian means and covariances at that time, render the image $\hat I_t$, and optimize
\begin{equation}
\mathcal{L}
=
\mathcal{L}_{\text{rec}}
+
\lambda_f \mathcal{L}_{\text{flow}}
+
\lambda_r \mathcal{L}_{\text{ode}}.
\end{equation}

\subsection{Reconstruction Loss}

The reconstruction term combines pixelwise fidelity and structural similarity:
\begin{equation}
\mathcal{L}_{\text{rec}}
=
\sum_t
\|\hat I_t-I_t\|_1
+
\lambda_s
\sum_t
\left(
1-\operatorname{SSIM}(\hat I_t,I_t)
\right).
\end{equation}

\subsection{Optical Flow Consistency Loss}

To encourage temporally coherent motion, we compare the optical flow between consecutive rendered frames and the optical flow between consecutive ground-truth frames:
\begin{equation}
\mathcal{L}_{\text{flow}}
=
\sum_t
\left\|
\hat F_{t\to t+1}
-
F_{t\to t+1}
\right\|_1.
\end{equation}
Here, $F_{t\to t+1}$ denotes the optical flow estimated from the ground-truth images $I_t$ and $I_{t+1}$, while $\hat F_{t\to t+1}$ denotes the optical flow estimated from the rendered images $\hat I_t$ and $\hat I_{t+1}$.

\subsection{ODE Regularization}

To prevent unstable trajectories and degenerate covariance evolution, we regularize the ODE dynamics:
\begin{equation}
\mathcal{L}_{\text{ode}}
=
\mathcal{L}_{\text{traj}}
+
\lambda_{\omega}\mathcal{L}_{\omega}
+
\lambda_s\mathcal{L}_{s}.
\end{equation}

The trajectory regularizer penalizes excessive 3D Gaussian velocity:
\begin{equation}
\mathcal{L}_{\text{traj}}
=
\sum_i
\int
\left\|
\frac{d x_i(\tau)}{d\tau}
\right\|_2^2
d\tau.
\end{equation}

Using the mean ODE, this becomes
\begin{equation}
\mathcal{L}_{\text{traj}}
=
\sum_i
\int
\left\|
A_i x_i(\tau)+b_i
\right\|_2^2
d\tau.
\end{equation}

The covariance orientation regularizer is
\begin{equation}
\mathcal{L}_{\omega}
=
\sum_i
\left\|
\omega_i^\Sigma
\right\|_2^2.
\end{equation}

The covariance scale regularizer is
\begin{equation}
\mathcal{L}_{s}
=
\sum_i
\int
\sum_{j=1}^3
\left[
\left(
\frac{d s_{ij}(\tau)}{d\tau}
\right)^2
+
\beta |s_{ij}(\tau)|
\right]
d\tau.
\end{equation}

Since
\begin{equation}
\frac{d s_{ij}(\tau)}{d\tau}
=
\kappa_{ij},
\end{equation}
this penalizes excessive scale growth or shrinkage.

\section{Training Procedure}

For a training frame at time $\tau$, we begin from the canonical Gaussian state at $\tau=0$. We solve the explicit real-valued ODEs for the mean and covariance to obtain
\begin{equation}
\mu_i(\tau),
\qquad
\Sigma_i(\tau)
\end{equation}
for all Gaussians. These updated Gaussian parameters are then passed directly into the original 3DGS renderer.

The reconstruction and temporal consistency losses backpropagate through the renderer and through the ODE solution, thereby updating the canonical Gaussian parameters and the ODE parameters controlling mean and covariance evolution. As a result, the model learns continuous-time 3D trajectories for Gaussian position and shape while remaining fully compatible with standard 3DGS.

\section{Evaluation Protocol}

Since the main objective of this work is to evaluate whether explicit ODE-based Gaussian evolution improves deformable motion modeling, evaluation should not rely only on conventional framewise image metrics. We therefore use two types of evaluation.

\subsection{Framewise Rendering Metrics}

As standard rendering-quality measures, we report image-based metrics such as PSNR, SSIM, and LPIPS between rendered frames and ground-truth frames. These measure per-frame reconstruction fidelity.

\subsection{Motion-Oriented Evaluation and Visualization}

To better analyze deformable motion behavior over time, we additionally use optical-flow-based and trajectory-based observations.

First, we measure optical flow error by comparing the optical flow estimated from consecutive rendered frames with the optical flow estimated from consecutive ground-truth frames. This provides a motion-oriented metric for whether the rendered temporal evolution matches the true image-space motion.

Second, we analyze trajectory smoothness and jitter from the learned Gaussian trajectories themselves. In particular, we visualize Gaussian center trajectories, velocity curves, acceleration curves, and jitter heatmaps over time.

Therefore, optical flow error and trajectory smoothness are used not only as evaluation criteria but also as visualization tools for understanding how well the temporal model captures deformable motion.

\section{Comparison Setting and Baselines}

The framework is designed to evaluate temporal modeling strategies under a fixed rendering backbone. Since rendering is always performed by the original 3DGS pipeline, differences in performance arise only from how Gaussian mean and covariance evolve over time.

The main methods to compare are:
\begin{itemize}[leftmargin=1.5em]
    \item \textbf{3D real-valued ODE:} Gaussian mean evolves directly by an explicit real-valued ODE in $\mathbb{R}^3$, and covariance evolves by explicit ODEs for scale and orientation.
    \item \textbf{Linear approximation:} Gaussian mean and covariance are modeled as linear functions of time.
    \item \textbf{Deformation network:} Gaussian mean and covariance are updated by a learned deformation module.
    \item \textbf{Framewise optimization baseline:} Gaussian parameters are optimized independently or only weakly coupled across frames.
\end{itemize}

This setup directly evaluates whether explicit real-valued ODE-based temporal evolution can track deformable motion more effectively than standard alternatives.

\section{Discussion}

The proposed method intentionally keeps the original 3DGS formulation intact. The scene is still represented by explicit 3D Gaussians, and rendering is still performed by projection, splatting, and alpha compositing. The contribution lies only in the temporal parameterization of Gaussian geometry.

By modeling Gaussian mean and covariance as explicit real-valued ODE-driven quantities, the framework introduces a compact and interpretable mechanism for dynamic motion. The key difference from planar or complex-valued formulations is that the mean trajectory is modeled directly in $\mathbb{R}^3$. Therefore, the motion equation is naturally aligned with the geometry of the 3D Gaussian primitive.

This avoids the need to define a local 2D motion plane and avoids representing 3D motion through a complex-valued planar state. Instead, each Gaussian has a unified 3D dynamical system controlling its center, while its shape evolves through explicit scale and rotation ODEs.

Because no warp field or implicit neural deformation is inserted into the renderer, the experiments offer a clean test of whether explicit real-valued ODE dynamics alone are sufficient to improve deformable motion tracking in dynamic 3DGS.

\section{Conclusion}

We presented a dynamic 3D Gaussian Splatting framework in which the original 3DGS renderer is preserved and only the temporal evolution of Gaussian mean and covariance is changed. The Gaussian mean is modeled directly by a real-valued ODE in three-dimensional space, while the covariance is evolved through explicit ODEs for scale and orientation. This creates a unified and natural 3D formulation for continuous-time Gaussian motion, allowing the framework to evaluate whether explicit ODE dynamics can better represent deformable motion than linear approximation or deformation-network-based approaches, while remaining faithful to the original 3DGS pipeline.

\end{document}